\cotefile{138}
\scanpage{60}
\deleted{$C=(C_0,C_1)$ \unsure{groupoide|groupoïde} connexe

$G=\operatorname{Aut}(C)$ groupe des automorphismes

$G^0$ ss-groupe des $g\in G$ induisant l'identité sur $C_0$

\diagram{\begin{tikzcd}[ampersand replacement=\&,column sep=small]
1\arrow[r] \& G^0\arrow[r] \& G\arrow[r,swap,"\text{\uuline{surj}}"] \& \unsure{\mathfrak S|\mathscr S}_{C_0}\arrow[r] \& 1
\end{tikzcd}}

[ Pour qu'une \unsure{autoéquivalence} \inserted{$u$} soit un iso, il f et s. qu'elle soit bijective sur les objets de $C$, et induise un iso. \deleted{sur les} $\operatorname{Aut}(a)\longrightarrow\operatorname{Aut}(u(x))$ pour \uline{un} $x\in C_0$ ; d'autre part, pour toute application $C_0\overset{u_0}{\longrightarrow}C'_0$}

\diagram{\begin{tikzcd}[ampersand replacement=\&,column sep=small]
1\arrow[r,dash] \& \operatorname{Aut}^0(C)\arrow[r,dash] \& \operatorname{Aut} C\arrow[r,dash] \& \mathfrak S_C\arrow[r,dash] \& 1
\end{tikzcd}}

\[
\deleted{\operatorname{Aut}^0(C)\simeq}
\]
\diagram{\begin{tikzcd}[ampersand replacement=\&,column sep=small]
\unsure{0|o}\mathord{-}\cdot \\
1\arrow[r] \& \operatorname{Aut}^{\unsure{0}}(C,a)\arrow[r,dash] \& \operatorname{Aut}^0(C)\arrow[r,swap,"\unsure{1}"] \& \operatorname{Aut}\bigl(\underset{\substack{\Vert\\\operatorname{Aut}_C(a)}}{\pi_a}\bigr)\arrow[r] \& 1
\end{tikzcd}}

\deleted{A} Pour déterminer $\deleted{\operatorname{Aut}^0(C,a)=}$
\[
\operatorname{Aut}^0(C,a)=\bigl\{f\in\operatorname{Aut}(C)\mid \operatorname{Ob} f=\mathrm{id}\ :\bigl(f_a:\operatorname{Aut}(a)\longrightarrow\operatorname{Aut}(\underset{\substack{\Vert\\f(a)}}{a})\bigr)=\mathrm{id}\bigr\}
\]
On considère, pour tt $x\in C_0$,
\[
T_x=\operatorname{Isom}(a,x)\ ,\ \pi_a\text{-torseur à droite}
\]
et son commutant $\deleted{\unsure{G_x}}\operatorname{Aut}(x)$
% archive sheet number: 59
